% basic nastaveni
\documentclass[12pt]{article}
\setlength{\parindent}{0pt}
\setlength{\parskip}{1em}
\usepackage[utf8]{inputenc}
\usepackage[czech]{babel}
\usepackage[T1]{fontenc}
\usepackage{graphics}
\usepackage{caption}
\usepackage{hyperref}
\usepackage{dsfont}
\usepackage{color}
\usepackage{float}
\usepackage{enumitem}
\usepackage{graphicx}

% matematicke package
\usepackage{tikz-cd}
\usepackage{tikz}
\usetikzlibrary{babel}
\usepackage{amsmath}
\usepackage{amssymb}
\usepackage{amsfonts}
\usepackage{amssymb}
\usepackage{geometry}
\geometry{margin=2.5cm}
\usepackage{pgfplots}
\pgfplotsset{compat=1.18}
\usepackage{mathrsfs}
\usetikzlibrary{shapes.misc}
\usetikzlibrary{angles,quotes}

\title{Maturitní otázka číslo 8: Nealgebraické rovnice, nerovnice a jejich soustavy}
\author{}
\date{}

\begin{document}
\maketitle

\section*{Exponenciální rovnice a nerovnice}
Základem je vztah $a^{f(x)} = a^{g(x)}$, kde $a> 0, a\neq 1$.

\textbf{Metody řešení}
\begin{itemize}[itemsep = 2pt]
    \item Úprava na stejný základ ($a$): Pokud lze obě strany převést na mocniu o stejném základu, využíváme \textbf{prostoty exponenciální funkce}: $a^{f(x)} = a^{g(x)} \Leftrightarrow f(x) = g(x)$
    \item Substituce: Typicky u rovnic tvaru $A \cdot a^{2x} + B\cdot a^x + C = 0$. Zavedeme substituci $t = a^x$ ($t$ bude vždy větší než nula) a řešíme kvadratickou rovnici.
    \item Logaritmování: Pokud základy nelze sjednotit (například $7^x = 2^{x-1}$), logaritmujeme obě strany rovnice
    
    Příklad: řešte rovnici $7^x = 2^{x-1}$

    Řešení:
    \begin{align*}
        \log_7 (7^x) &= \log_7 (2^{x-1})\\
        x \cdot \log_7 7 &= (x-1) \cdot \log_7 2\\
        x\cdot (1 - \log_7 2) &= -\log_7 2\\
        x &= \frac{-\log_7 2}{1 - \log_7 2}\\
    \end{align*}
    
\end{itemize}

\subsection*{Exponenciální nerovnice}
Při odstraňování základů u nerovnic $a^{f(x)} > a^{g(x)}$je důležité sledovat hodnotu $a$

Pokud $a>1$, tak je funkce rostoucí a znaménko nerovnosti se nemění

Pokud je $a<1$, tak je funkce klesající a znaménko nerovnosti se obrací

\section*{Logaritmické rovnice a nerovnice}
Logaritmus $log_a x$ je definován pouze pro $x>0$ a $a \in (0,1) \cup (1,\infty)$, takže nezapomeňte na podmínky

\textbf{Věty o logaritmech}
\begin{itemize} [itemsep = 2pt]
    \item $\log_a x = \log_a y = \log_a(x\cdot y)$
    \item $\log_a x - \log_a y = \log_a(\frac{x}{y})$
    \item $\log_a x^k = k \cdot \log_a x$
\end{itemize}

Odlogaritmování můžeme použít, pokud převedeme rovnici do tvaru $\log_a f(x) = \log_a g(x)$, díky prostosti logaritmických funkcí, tak můžeme rovnici upravit na $f(x) = g(x)$

\subsection*{Logaritmické nerovnice}
Pro $a>1$ platí $\log_a f(x) < \log_a g(x) \Leftrightarrow f(x) < g(x)$

Pro $0<a<1$ platí $\log_a f(x) < \log_a g(x) \Leftrightarrow f(x) > g(x)$

A výsledný interval musí být v průniku s definičním oborem nerovnice.

\section*{Goniometrické rovnice a nerovnice}
\textbf{Jednotková kružnice} je kružnice se středem v počátku souřadnic o poloměru 1.

Základní rovnice jako $\sin x = a$ řešíme pomocí jednotkové kružnice nebo grafu, typicky existuje více řešení.

Složitější rovnice, jako $2\sin^2 x + 3\sin x + 1 = 0$ řešíme pomocí substituce. Tento konkrétní příklad subsituujeme jako $t = \sin x, |t|\leq 1$

\textbf{Zápis řešení}

Typicky dostaneme několik základních řešení v rámci jedné periody (fundamentální řešení), ke kterým přičítáme násobky periody.

Obecně množinu řešení zapisujeme jako množinu $K$ a vypadá následovně:
\[K = \underset{i=1}{\overset{n}{\bigcup}}\{x_i + k\cdot T; k\in\mathbb{Z}\}\]
Kdokoliv se tohle učí, tak je to  sjednocení zapsaný blbě, brzo to opravim (takhle ho po nás ve škole nechtěj)
\end{document}
